\subsection{常规去噪基线方法}
\label{sec:conventional-baselines}

经典去噪方法构成了激光超声信号增强的标准基线。Wiener滤波通过线性最小均方误差(MMSE)估计信号，假设平稳统计量；
然而，其频域平滑破坏了携带缺陷指示信息的高频瞬态分量。基于小波的去噪(DWT)将信号分解为多分辨率系数，并应用软阈值或硬阈值来抑制噪声\cite{chen2019wavelet}。\texttt{db4}小波因其良好的时频局部化特性常被用于激光超声信号\cite{chen2019wavelet}。BM3D通过将相似信号块分组为3-D块并应用协同滤波来扩展小波阈值，在聚合噪声指标上实现最先进的性能\cite{dabov2007bm3d}。

这些方法的共同根本局限：所有方法都优化相同的目标函数——MSE或SNR改进——而没有对声学特征保留的显式约束。
对于激光超声缺陷检测，诊断相关特征包括Lamb波模态的到达时间、模态间的振幅比以及波形形态。
这些特征在去噪过程中未受到显式惩罚，而激进的噪声抑制通常会降低它们。陈等人指出，小波阈值即使在SNR改善时也可能扭曲到达时间\cite{chen2019wavelet}，且BM3D基于块的处理引入额外的相位失真。

因此，没有常规方法显式保留激光超声缺陷检测的特征。这一差距促使将声学特征保留作为一等优化目标的物理约束方法。
